

\ExplSyntaxOn

\prop_new:N \l_zig_node_content_prop
\prop_new:N \l_zig_arrow_label_prop

% Expand the key before looking it up.
\cs_generate_variant:Nn \prop_item:Nn { Ne }
\cs_generate_variant:Nn \prop_if_in:NnTF { Ne }

% Look up node content; use the fallback if there is no entry.
\NewDocumentCommand{\ObjectContent}{m m}
 {
  \prop_if_in:NeTF \l_zig_node_content_prop {#1}
   { \prop_item:Ne \l_zig_node_content_prop {#1} }
   { #2 }
 }

% Look up morphism label; use the fallback if there is no entry.
\NewDocumentCommand{\MorphismContent}{m m}
 {
  \prop_if_in:NeTF \l_zig_arrow_label_prop {#1}
   { \prop_item:Ne \l_zig_arrow_label_prop {#1} }
   { #2 }
 }

\NewDocumentCommand{\SetObjectContent}{m m}{
  \prop_put:Nen \l_zig_node_content_prop {#1} {#2}
}

\NewDocumentCommand{\SetMorphismContent}{m m}{
  \prop_put:Nen \l_zig_arrow_label_prop {#1} {#2}
}

\ExplSyntaxOff

%%%%%%%%%%%%%%%%%%%%%%%%%%%%%%%%%%%%%
% \ob and \Ob macros used in pstricks diagrams
%%%%%%%%%%%%%%%%%%%%%%%%%%%%%%%%%%%%%

% nodenames are constructed and have an optional instance
%  parameter so that two copies of the same content 
%  can be separately identified.

% Construct node names
\newcommand{\ob}[3]{#1#2#3}
\newcommand{\bob}[1]{#1}
\newcommand{\obp}[3]{#1#2#3p}
\newcommand{\obs}[3]{#1#2#3s}
\newcommand{\obpp}[3]{#1#2#3pp}
\newcommand{\obss}[3]{#1#2#3ss}

% Canonical node-name constructor:
% {instance}{object}{subscript}{operator}
\newcommand{\nodename}[4]{#1#2#3#4}

% Object
\NewDocumentCommand{\Ob}{O{} m m}{%
  \Rnode{\nodename{#1}{#2}{#3}{}}{#2_{#3}}%
}

% Bare object, without subscript
\NewDocumentCommand{\bOb}{O{} m}{%
  \Rnode{\nodename{#1}{#2}{}{}}{#2}%
}

% Predecessor
\NewDocumentCommand{\Obp}{O{} m m}{%
  \Rnode{\nodename{#1}{#2}{#3}{p}}{#2_{#3-1}}%
}

% Successor
\NewDocumentCommand{\Obs}{O{} m m}{%
  \Rnode{\nodename{#1}{#2}{#3}{s}}{#2_{#3+1}}%
}

% Two predecessors
\NewDocumentCommand{\Obpp}{O{} m m}{%
  \Rnode{\nodename{#1}{#2}{#3}{pp}}{#2_{#3-2}}%
}

% Two successors
\NewDocumentCommand{\Obss}{O{} m m}{%
  \Rnode{\nodename{#1}{#2}{#3}{ss}}{#2_{#3+2}}%
}


% \zigzagTwo[instance]{x}{y}{f}
% draws a zigzag of length two <f1,f'1,f2> 
% where f1 : x1 -> y1
%       f'1: x2 -> y1
%       f2 : x2 -> y2
%
% overide xi  using \SetObjectContent{\ob{instance}{x}{i}}{...}
% override fi using \SetMorphismContent{fi}{...}
%                or \SetMorphismContent{fprimei}{...} 


\NewDocumentCommand{\zigzagTwoObjects}{O{} m m}{
  \begin{array}{c c c c}
   &\Rnode{\ob{#1}{#3}{1}}
      {\ObjectContent{\ob{#1}{#3}{1}}{#3_{1}}}
   &&
     \Rnode{\ob{#1}{#3}{2}}
      {\ObjectContent{\ob{#1}{#3}{2}}{#3_{2}}}
   \\[0.8cm]
   \Rnode{\ob{#1}{#2}{1}}
      {\ObjectContent{\ob{#1}{#2}{1}}{#2_{1}}}
   &&
   \Rnode{\ob{#1}{#2}{2}}
      {\ObjectContent{\ob{#1}{#2}{2}}{#2_{2}}}
  \end{array}
}
\NewDocumentCommand{\zigzagTwoObjectsBulleted}{O{} m m}{
  \begin{array}{c c c c}
   &\Rnode{\ob{#1}{#3}{1}}
      {\ObjectContent{\ob{#1}{#3}{1}}{\bullet}}
   &&
     \Rnode{\ob{#1}{#3}{2}}
      {\ObjectContent{\ob{#1}{#3}{2}}{\bullet}}
   \\[0.8cm]
   \Rnode{\ob{#1}{#2}{1}}
      {\ObjectContent{\ob{#1}{#2}{1}}{\bullet}}
   &&
   \Rnode{\ob{#1}{#2}{2}}
      {\ObjectContent{\ob{#1}{#2}{2}}{\bullet}}
  \end{array}
}
\NewDocumentCommand{\zigzagTwoMorphisms}{O{} m m m}{
  \begin{arrows}
    \ncarr{\ob{#1}{#2}{1}}{\ob{#1}{#3}{1}}
      \alabel{
        \MorphismContent{#1#4one}{#4_{1}}
      }[0.6][0.4]
    \ncarr{\ob{#1}{#2}{2}}{\ob{#1}{#3}{1}}
      \blabel{
        \MorphismContent{#1#4primeone}{#4'_{1}}
      }[0.6][0.4]
    \ncarr{\ob{#1}{#2}{2}}{\ob{#1}{#3}{2}}
      \blabel{
        \MorphismContent{#1#4two}{#4_{2}}
      }[0.7][0.4]
  \end{arrows}
}

\NewDocumentCommand{\zigzagTwo}{O{} m m m}{
  \zigzagTwoObjects[#1]{#2}{#3}
  \zigzagTwoMorphisms[#1]{#2}{#3}{#4}
}
\NewDocumentCommand{\zigzagTwoBulleted}{O{} m m m}{
  \zigzagTwoObjectsBulleted[#1]{#2}{#3}
  \zigzagTwoMorphisms[#1]{#2}{#3}{#4}
}

